% ============================================================
% QUESTION TYPE: FILL IN THE BLANKS
%
% PEANUT EXAM FILL-IN-THE-BLANKS LAYOUT STANDARD
%
% AI INSTRUCTION:
% - Follow this format EXACTLY for fill-in-the-blank questions
% - Questions are bold-numbered
% - Blanks are created using \underline{\hspace{<width>}}
% - Extra context or hints may appear in parentheses
% - Use \vspace{...} for spacing only
% - Do NOT convert to multiple choice or written format
% - Do NOT add answer boxes
% ============================================================

\section*{Part II: [PART TITLE]}
\textit{Fill in the blanks with the most appropriate answer.
Extra context or instructions may be provided in parentheses.}

\vspace{2mm}

\textbf{1.}
This is a sentence with a blank:
The system \underline{\hspace{3cm}} (to operate) continuously for several hours.
(It is still running now.)

\vspace{4mm}

\textbf{2.}
Another sentence testing grammar, terminology, or usage:
The report \underline{\hspace{3cm}} (to submit) yesterday. (Use passive voice.)

\vspace{4mm}

\textbf{3.}
This question may involve tense, agreement, or modality:
They \underline{\hspace{3cm}} (not / to complete) the task yet.
(Make it negative.)

\vspace{4mm}

\textbf{4.}
By the time the experiment started, the data
\underline{\hspace{3cm}} (to collect).

\vspace{4mm}

\textbf{5.}
Technical documents must \underline{\hspace{3cm}} clearly to avoid ambiguity.
(Use modal verb.)

\vspace{6mm}
